\documentclass[11pt]{article}
\usepackage{lmodern}          % AER font
\usepackage[T1]{fontenc}      % Better font encoding
\usepackage{amsmath, amssymb}
\usepackage[margin=1in]{geometry}
\usepackage{graphicx}
\usepackage{url}
\usepackage[round, longnamesfirst]{natbib}
\usepackage{bibentry}         % lets \nobibliography resolve cites without printing a list
\usepackage{hyperref}

%=======================================
% Format helpers
%=======================================
% Target school -- used in the header and the conclusion so they stay in sync.
\newcommand{\school}{[University name]}
\newcommand{\program}{[Program name]}

% Run-in section heading: bold label ending in a period, body text follows
% on the same line (matches the reference format).
\newcommand{\shead}[1]{\par\vspace{0.85em}\noindent\textbf{#1.}\ }

% Indented paragraphs, no extra vertical space between them.
\setlength{\parindent}{1.5em}
\setlength{\parskip}{0pt}

% Word count footer. The VS Code build runs texcount first, writing the count
% to \jobname.wc; it shows "??" until texcount has run once.
\usepackage{fancyhdr}
\newread\wcfile
\IfFileExists{\jobname.wc}{\openin\wcfile=\jobname.wc \read\wcfile to\wcnum \closein\wcfile}{\def\wcnum{??}}
\pagestyle{fancy}\fancyhf{}\renewcommand{\headrulewidth}{0pt}
\fancyfoot[C]{\thepage}\fancyfoot[R]{\small Word count: \wcnum}

\begin{document}

%=======================================
% Header (not counted in the word count)
%=======================================
%TC:ignore
\noindent
\begin{minipage}[t]{0.60\textwidth}
  {\large\bfseries Statement of Purpose}
\end{minipage}\hfill
\begin{minipage}[t]{0.38\textwidth}
  \raggedleft
  Rhea Choudhary\\[2pt]
  \school
\end{minipage}

\vspace{0.5em}
\noindent\rule{\textwidth}{0.8pt}
%TC:endignore

\vspace{1.1em}

%=======================================
% Body
% 500-word version of Rhea_SOP.tex. Edits to the full version do NOT flow
% through here -- update both.
%=======================================
\shead{Research interests} I am applying to pursue a PhD in \program{} at \school, to study how the design of social programs affects labour market outcomes.

\shead{Academic background} I completed a Bachelor of Commerce (Honours) majoring in Economics and Finance at the University of Melbourne, achieving First Class Honours and numerous academic awards. My Honours thesis, on unconventional monetary policy and financial stability in an emerging market, was published as a central bank working paper. I also excelled in Real Analysis and Linear Algebra.

\shead{Pre-doctoral and professional experience} After working as a senior analyst at the Reserve Bank of Australia, I now work as a predoctoral researcher for Associate Professor Manasi Deshpande. Our current project exploits the staggered rollout of the NDIS, a first-of-its-kind program funding care services for people with disabilities, to estimate its effects on wages and labour supply. We find that the marginal entrants into disability work are low-income, lower-skilled women with modest earnings gains, while skilled occupations such as therapists see higher wages and lower entry. This has given me extensive experience with large administrative datasets and causal inference methods such as event studies and shift-share instruments.

\shead{Current research} I am interested in the relative merits of different approaches to increase labour supply in high-demand sectors. A common mechanism is raising wages (e.g., \cite{dalbo_finan_rossi2013}), but its employment effect is ambiguous in equilibrium, since it also raises firm labour costs. And if reskilling is costly, the wage needed to compensate may be too high. Alternatively, subsidising education reduces fixed costs of entry \citep{schoonover2026}. Which works better depends on the labour supply elasticity of earnings versus entry costs. \\~\\I study this using two Australian policies. Since 2019, Victoria's `Free TAFE' has offered free education in courses for priority occupations, with residency as the main eligibility condition. A national 15 per cent wage increase for aged care workers from 2023 \citep{fwc_agedcare2022} is similarly sector-targeted, and because the federal government funds it, firms have less incentive to change their input mix. \\~\\Starting with the education subsidy, I use Australian administrative data and a difference-in-differences comparing Victoria to other states. I find that Free TAFE generated a large and immediate increase in enrolments, particularly among women, working-aged cohorts and in more advanced courses. Constructing a `risk set' of likely enrollees from ex-ante characteristics, I find that young women previously on income support have the strongest enrolment first stage, which I verify using a distance-to-college instrument \citep{card1995}. I am now estimating employment and wage responses, to compare against entry responses to the aged care increase.

\shead{Future research} This work sits within a broader agenda on how social program design, especially support beyond income assistance, shapes labour and welfare outcomes. Family planning, for instance, also affects labour supply (\cite{gallen2026}). I intend to examine Victoria's `free IVF' initiative, which offers IVF at public fertility clinics with no means test. By lowering a fixed cost of having a child later rather than replacing lost income, it could shape labour market and childbirth decisions differently from parental income assistance.

% Citations resolve against bib.bib, but no reference list is printed.
\bibliographystyle{abbrvnat}
\nobibliography{bib}

\end{document}
